\chapter{Muscle Atrophy}

\section*{Definition}
Muscle atrophy is the progressive loss of skeletal muscle mass and strength, characterized by a reduction in cross-sectional area of myofibers. It is classified as disuse (non-neurogenic) or neurogenic, depending on whether lower motor neuron innervation is intact.

\section*{What Happens in Your Body}
Skeletal muscle homeostasis depends on a balance between protein synthesis and degradation. Atrophy results from increased ubiquitin-proteasome pathway activity, autophagy, and caspase-mediated apoptosis. Disuse atrophy involves decreased IGF-1/Akt signaling and upregulation of myostatin. Neurogenic atrophy follows denervation, leading to spontaneous depolarizations (fibrillations) and grouped fiber-type loss. Chronic inflammation elevates TNF-$\alpha$ and IL-6, further driving catabolic signaling.

\section*{Causes}
\begin{itemize}
  \item \textbf{Disuse:} Prolonged immobilization, bed rest, sedentary lifestyle, casting
  \item \textbf{Neurogenic:} Peripheral neuropathy (diabetic, alcoholic), radiculopathy, spinal cord injury, amyotrophic lateral sclerosis (ALS), polio, Guillain-Barr\'{e} syndrome
  \item \textbf{Myopathic:} Muscular dystrophies (Duchenne, Becker), inflammatory myopathies (polymyositis, dermatomyositis, inclusion body myositis)
  \item \textbf{Cachexia:} Malignancy, chronic heart failure, COPD, chronic kidney disease, HIV/AIDS
  \item \textbf{Endocrine:} Cushing syndrome, hyperthyroidism, hypogonadism, growth hormone deficiency
  \item \textbf{Medication-related:} Glucocorticoids, statins, colchicine, antiretrovirals
  \item \textbf{Metabolic:} Starvation, malabsorption, vitamin D or E deficiency
\end{itemize}

\section*{When to See a Doctor}
See your doctor if:
\begin{itemize}
  \item Rapidly progressive weakness or muscle loss over weeks to months
  \item Asymmetric atrophy or focal muscle wasting
  \item Loss of ambulation or increasing falls
  \item Dysphagia, dysarthria, or respiratory difficulty
\end{itemize}

\section*{Related Symptoms}
\begin{itemize}
  \item Progressive weakness, often proximal (rising from chair, climbing stairs)
  \item Fasciculations (suggest neurogenic cause)
  \item Loss of deep tendon reflexes in neuropathic atrophy
  \item Cramping and myalgias
  \item Gower sign in proximal myopathy (using hands to push off floor when rising)
\end{itemize}
\textbf{Important:} Fasciculations with atrophy strongly suggest lower motor neuron disease (e.g., ALS) and warrant neurological evaluation with electromyography.

\section*{Self-Care / Home Management}
\begin{itemize}
  \item Range-of-motion exercises to prevent contractures in immobilized limbs
  \item Nutritional optimization with adequate protein intake (1.2--1.5 g/kg/day)
  \item Hydration and electrolyte balance to support muscle function
  \item Fall prevention measures: remove hazards, use assistive devices
  \item Avoidance of prolonged bed rest or immobilization when possible
\end{itemize}

\section*{How Your Doctor Will Evaluate You}
\begin{itemize}
  \item Electromyography and nerve conduction studies to distinguish neurogenic vs. myopathic causes
  \item Muscle biopsy (open or needle) for histopathology and enzyme analysis
  \item Laboratory evaluation: CK, aldolase, TSH, cortisol, vitamin B12, vitamin D, ANA, ESR
  \item Genetic testing for dystrophinopathies and other inherited myopathies
  \item MRI of affected muscles for fatty infiltration and inflammation
  \item Nerve imaging (ultrasound or MRI) for structural neuropathies
\end{itemize}

\section*{Treatment Options}
\begin{itemize}
  \item Physical therapy with progressive resistance and neuromuscular electrical stimulation
  \item Discontinue causative medications where possible
  \item Corticosteroids or immunosuppressants for inflammatory myopathies
  \item Enzyme replacement or gene therapy for certain dystrophies (e.g., eteplirsen for Duchenne)
  \item Nutritional support and appetite stimulants for cachexia
  \item Disease-modifying therapies for ALS (riluzole, edaravone) and other neurodegenerative disorders
  \item Orthotic devices and mobility aids to maintain function
\end{itemize}

\section*{References}

\begin{thebibliography}{99}
\bibitem{harrison-muscle} Longo DL, et al. \emph{Harrison's Principles of Internal Medicine}. 21st ed. McGraw-Hill; 2022. Chapter 436: Muscle Atrophy.
\bibitem{uptodate-atrophy} Tawil R, et al. Approach to the patient with muscle atrophy. In: UpToDate, Post TW (Ed), UpToDate, Waltham, MA; 2025.
\bibitem{cohen} Cohen S, et al. Muscle wasting: molecular mechanisms and therapeutic targets. \emph{Nat Rev Drug Discov}. 2023;22(6):471-490.
\bibitem{evans} Evans WJ, et al. Sarcopenia: a clinical practice guideline. \emph{J Am Med Dir Assoc}. 2024;25(1):12-23.
\bibitem{scarano} Scarano M, et al. Neurogenic atrophy and motor neuron disease. \emph{Neurology}. 2023;100(12):567-578.
\bibitem{keda} Ikeda S, et al. \emph{Neuromuscular Disorders}. 4th ed. Elsevier; 2024. Chapter 8: Atrophic Disorders.
\end{thebibliography}

\clearpage
